\documentclass[10pt,a4paper]{amsart}
\usepackage[margin=19mm]{geometry}
\usepackage{amsmath,amssymb,mathtools}
\usepackage[scr=rsfs,cal=euler]{mathalfa}
\usepackage{xfrac}
\usepackage[unicode,hidelinks,bookmarks=false]{hyperref}
\newcommand{\ie}{i.e.\ }
\newcommand{\cf}{cf.\ }
\newcommand{\resp}{respectively\ }
\newcommand{\Subsection}{\subsection}
\providecommand{\nobreakdash}{\nobreak\hbox{-}}
\setlength{\emergencystretch}{2em}
\multlinegap0pt
\newtheorem{prop}{Proposition}
\newcommand{\De}{\Delta}
\newcommand{\Om}{\Omega}
\newcommand{\bt}{\mathbf{t}}
\newcommand{\dd}{\partial}
\newcommand{\de}{\delta}
\newcommand{\la}{\lambda}
\newcommand{\na}{\nabla}
\newcommand{\om}{\omega}
\newcommand{\sC}{{\mathscr C}}
\newcommand{\sE}{{\mathscr E}}
\newcommand{\sH}{{\mathscr H}}
\newcommand{\sU}{{\mathscr U}}
\newcommand{\sg}{\sigma}
\title{$G_2$-Laplacian flow on complete asymptotically conical manifolds}
\author{Ilyas Khan}
\date{Source: arXiv:2608.26042v1 (CC BY 4.0)}
\begin{document}
\maketitle

\noindent\emph{Source and licence.} $G_2$-Laplacian flow on complete asymptotically conical manifolds. Version arXiv:2608.26042v1 (CC BY 4.0), distributed by arXiv under the Creative Commons Attribution 4.0 International licence, http://creativecommons.org/licenses/by/4.0/, as recorded in the arXiv licence metadata for this exact version and shown on the abstract page https://arxiv.org/abs/2608.26042v1. Full licence notice: \texttt{licenses/arxiv-2608.26042-CC-BY-4.0.txt}.\par\smallskip

\noindent\textit{Editorial scope.} Counted unit: the article's prop forms (source0.tex lines 2894--2920). The article's complete original proof of it is delivered with it (source0.tex lines 2922--2947, one \texttt{proof} environment of 26 lines). No mathematics is written, repaired or re-derived: the statement and the proof are the article's own lines, in the article's own notation, and no retained line is substituted or re-flowed.  Retained uncounted as the source-local context the proof needs: lines 749--751; lines 2039--2041; lines 2135--2175.  Nothing was omitted from any retained span, so the package carries no omission record.  No retained line contains a citation, so the wrapper carries no bibliography and no bibliography entry.  No retained line contains a figure environment, an image-inclusion command or drawing code, so no image is delivered and the package declares no asset dependency.  Editorial formatting only, in addition: an AMS \texttt{amsart} preamble that restates the article's own declarations and macros used by the retained lines, namely the environments \texttt{prop} on separate counters, together with the standard packages those lines need.  The retained environments print in this document as Proposition 1 in order of appearance, whereas the article prints the same items with the numbering of its own full document.  The complete native source of the article as distributed in the arXiv e-print for this exact version is bundled unchanged as \texttt{sources/208616/source0.tex}.
\begin{multline}\label{equation helmholtz decomp}
\om_{cl} \in W^{s,p}\sE^k(N) \oplus W^{s,p}\sH_{D}^k(N) \text{ and }\bt \om_{cl} = 0,\\ \om_{co} \in W^{s,p}\sC^k(N) \oplus W^{s,p}\sH_{co}^k(N).
\end{multline}

\begin{equation}\label{equation conditions on coeffs}
|g^{ij}G^a_b - A^{aij}_b| = o(\rho^{-2}), |\na (g^{ij}G^a_b - A^{aij}_b)|, |B|, \text{ and } |C| = O(\rho^{-2}),
\end{equation}

\begin{prop}\label{proposition lopatinski shapiro condition}
Let $L$ be the elliptic operator of order $4$ defined on $k$-forms $\Omega^k(M)$ on a complete AC manifold $(M,g)$ by
\begin{align*}
    L\sg &= \De^*\De \sg \\
        &= \rho^{2\la + 4 + 7}\De (\rho^{-2\la -7} \De \sg) \\
        &= \rho^{4}\De \De \sg + C_3 \rho^{3} \na \rho * \na \De \sg + \rho^{2} (\De \rho )(\De \sg).
\end{align*}
Notice that the $j$th order coefficients $B_j$ have the property that
\begin{equation}\label{equation growth of coefficients}
|B_j| = O(1),
\end{equation}
for $j = 0, \ldots, 4$ (compare to the equivalent condition to \eqref{equation conditions on coeffs} in the parabolic case). Let $L^*$ denote the adjoint of $L$, defined by
\begin{align*}
    L^*\sg &= \De\De^* \sg \\
        &= \De \rho^{2\la + 4 + 7}\De (\rho^{-2\la -7} \sg) \\
        &= \rho^{4}\De \De \sg + \sum_{i=0}^3 \sum_{j=1}^{4-i} C_{ij} \rho^i \na^j \rho * \na^{4-i-j} \rho * \na^i \sg
\end{align*}
As before, the $j$th order coefficients $B_j$ have $O(1)$ decay. Let $M_R = \rho^{-1}((0,R]) \subset M$ and suppose that $\sg \in L^2_{\la}\Om^k(M_R) \cap H^4_{loc} \Om^k(M_R)$ solves either of the following adjoint boundary value problems: 
\begin{equation}\label{equation fourth order bvp with LS}
    \left\{
\begin{array}{lr}
L\sg = \eta   & \;\; \text{ on } M_R\\ 
\bt \sg = 0 \;\text{ and }\; \bt \rho \de \sg = 0 & \;\; \text{ on } \dd M_R,
\end{array}
\right.
\end{equation}
or its adjoint,
\begin{equation}\label{equation fourth order bvp with LS adjoint}
    \left\{
\begin{array}{lr}
L^*\sg = \eta   & \;\; \text{ on } M_R\\ 
\bt \sg = 0 \;\text{ and }\; \bt \rho^{2\la + 7 + 1} \de (\rho^{-2\la -7} \sg) = 0 & \;\; \text{ on } \dd M_R,
\end{array}
\right.
\end{equation}
where $\eta \in H_{\la}^{\ell-4}(M_R)$ for $\ell \ge 4$.
Then there exists a constant $c >0$ independent of $\sg, \eta,$ and $R$, such that
    \begin{equation}\label{equation schauder estimates forms with lopatinski shapiro}
        \|\sg\|_{H^{\ell}_{\la}(M_R)} \le c\Big(\|\eta\|_{H^{\ell-4}_{\la}(M_R)} + \|\sg\|_{L^2_{\la}(M_R)} \Big).
    \end{equation}
\end{prop}

\begin{prop}\label{proposition boundedness of restricted forms}
As before, let $U_R$ be the span of $\{\phi_i\}_{i= N}^\infty$, where $N$ is the smallest number such that for $i \ge N$,  $\sg_i^2 < C \|(\De^*\De)^{-1}\|$ for $C>0$ a fixed positive constant. Let $s \ge 3$ and let
\[
H_{\la}^{r-1,s-3}\sU^3(M_R \times I)
\]
be the image of $\mathrm{Proj}_U$ along time-slices of $H_{0,\la}^{r-1,s-3}\Om^3(M_R \times I)$ to $U_R$. Then, let 
\[
H_{0,\la}^{r-1,s-3}\sU^3(M_R \times I)
\]
be the image of the projection $\mathrm{Proj}_U$ along time-slices of the subspace $H_{0,\la}^{r-1,s-3}\Om^3(M_R \times I)$ of forms which vanish at the initial time to order $r-2$ in $t$. Then, if $\eta$ is in $H_{0,\la}^{r-1,s-3}\sU^3(M_R \times I)$, there is a potential $\xi \in H_{0,\la+1}^{r-1,s-2}\Omega_D^2(M_R \times I)$ such that
\[
d\xi = \mathrm{Proj}_{\sE} (\eta) \in H_{0,\la}^{r-1,s-3}\sE^3(M_R \times I),
\]
where $\mathrm{Proj}_{\sE}$ is the projection to the exact part of the Hodge-Morrey decomposition. Furthermore, there is a potential $\om \in H_{0,\la+1}^{r-1,s-2}\Omega^4(M_R \times I)$ such that 
\[
\de \om = \mathrm{Proj}_{\sC \oplus \sH_{co}} (\eta) \in H_{0,\la}^{r-1,s-3}(\sC^3 \oplus \sH_{co}^3) (M_R \times I),
\]
where $\mathrm{Proj}_{\sC \oplus \sH_{co}}$ is the projection to the co-exact piece of the Helmholtz decomposition given in \eqref{equation helmholtz decomp}. These two potentials satisfy the following estimates for all $t_0 \in I$. 
\begin{align}\label{equation timeslice potential estimates}
\|\dd_t^j \om (t_0)\|_{H_{\la+1 - 2j}^{s-2-2j}\Omega^4(M_R)},\; \|\dd_t^j \xi (t_0)\|_{H_{\la+1 - 2j}^{s-2-2j}\Omega^2(M_R)} \le C_M \|\dd_t^j \eta (t_0)\|_{H_{\la-2j}^{s-3-2j}\Omega^3(M_R)},
\end{align}
where $j \in \{0, \ldots, r-1\}$ and $C_M > 0$ does not depend on $R$. Note in particular that this implies that the maps 
\[
\mathrm{Proj}_{\sE} \text{ and } \mathrm{Proj}_{\sC \oplus \sH_{co}}
\]
are uniformly bounded maps on $H_{0,\la}^{r-1,s-3}\sU^3(M_R \times I)$.
\end{prop}

\begin{proof}
Let $\eta(t)$ be a family of 3-forms parametrized by $t$ in $H_{0,\la}^{r-1,s-3}\sU^3(M_R \times I)$. Fix $t \in I$. Then
\begin{align*}
    G(\eta(t)) &= \sum_i \sg_i \langle \phi_i, \eta(t) \rangle_{L^2_\la} \psi_i.
\end{align*}
This formula implies that the $L^2_{\la+2}$-norm of $G(\eta(t))$ is bounded by the $L^2_{\la}$-norm of $\eta(t)$ times a constant depending only on the operator norm $\|\De^{-1}(\De^{-1})^*\|$ on the whole space $M$ with respect to the initial metric $g_0$. This observation paired with the elliptic Schauder estimates in Proposition \ref{proposition lopatinski shapiro condition} gives higher order estimates in the spatial derivatives. Setting $\xi(t) := \de G( \eta(t))$ yields \eqref{equation timeslice potential estimates} when $j = 0$.

Since $\phi_i$ and $\psi_i$ are time-independent, differentiating the $\xi$ in time yields
\begin{align}\label{equation commute time and greens}
    \dd_t G(\eta) &= \dd_t \bigg(\sum_i \sg_i \langle \phi_i, \eta \rangle_{L^2_\la} \psi_i\bigg)\\
    &= \sum_i \sg_i \langle \phi_i, \dd_t \eta \rangle_{L^2_\la} \psi_i \nonumber\\
    &= G(\dd_t \eta) \nonumber
\end{align}
Identical calculations to those above yield the estimates in terms of the Sobolev norms of the time derivative $\dd_t \eta$. Performing these calculations for $j = 2,\ldots,r-1$ establishes that $G(\eta) \in H^{r-1, s-1}_{\la+2}\Omega^3_{hom} (M \times I)$. Since the metric is fixed at the initial time, the time derivative $\dd_t$ and the co-differential $\de$ commute. Set $\xi := \de G(\eta)$. Note that the homogeneous boundary condition on $G(\eta)$ implies that $\xi$ satisfies the Dirichlet boundary condition. Observe that for all $t_0 \in I$,
\begin{align*}
    \big\|\dd_t^j \xi \big|_{t_0} \big\|_{H_{\la+1 - 2j}^{s-2-2j}} &= \big\|\dd_t^j \de G(\eta) \big|_{t_0} \big\|_{H_{\la+1 - 2j}^{s-2-2j}} \\
    &= \big\|\de \dd_t^j G( \eta) \big|_{t_0} \big\|_{H_{\la+1 - 2j}^{s-2-2j}} \\
    &\le \big\|\dd_t^j G( \eta )\big|_{t_0} \big\|_{H_{\la+2 - 2j}^{s-1-2j}} \\
    &\le C_M \big\|\dd_t^j \eta \big|_{t_0} \big\|_{H_{\la-2j}^{s-3-2j}}.
\end{align*}
The last line follows from the commutativity of $G$ and $\dd_t$ shown in \eqref{equation commute time and greens}. In particular, this estimate implies that $\xi \in H^{r-1, s-2}_{\la+1}\Omega^2_{D} (M \times I)$. Furthermore, for $j = 0,\ldots, r-2$, this further implies that if $\eta \in H_{0,\la}^{r-1,s-3}\sU^3(M_R \times I)$
\[
\|\dd_t^j \xi (t_0)\|_{H_{\la+1 - 2j}^{s-2-2j}} \rightarrow 0, \; \text{ as } t_0 \rightarrow 0.
\]
Thus, $\xi \in H_{0,\la+1}^{r-1,s-2}\Omega_D^2(M_R \times I)$ and satisfies \eqref{equation timeslice potential estimates}. Setting $\om :=  dG(\eta)$ and repeating the proof given above with the natural modifications for the coexact case completes the proof of the theorem.
\end{proof}

\end{document}
